\scanpage{128}
\mspage{43}
% archive sheet number: 127

% L01-04
les $R^p p_*(R^qg'_*(\Lambda_\nu))|X_0\inserted{=R^p p_*(R^qg'_*(\Lambda_\nu)|X'_0),\quad\text{(th. de chgt de base propre)}}$, pour $\nu$ variable, forment un système AR-$\ell$-adique de faisceaux sur $X$, donc il en est de même des $R^i g_*(\Lambda_\nu)|X_0=(\underline{\Psi}_\nu^i)_{\mathrm{tame}}$, en vertu de (SGA 5 \uline{V} 5.2.1 et 5.2.4), d'où (i).

% L05-06
La dernière assertion de (ii) \inserted{relative à $(\underline{\Psi}^q)_{\mathrm{tame}}$ lui-même,} est une conséquence immédiate de ce qui précède.

% L07-10
\uline{Remarques} 6.3. Si $m'$ est l'entier défini dans 3.2 en termes des multiplicités des composantes de $X'_0$ dans $X'$, alors la démonstration précédente prouve que dans 6.2 on peut prendre pour $I'$ le s-groupe \inserted{ouvert,} d'indice $m'$ de $I$.

% L11-17
6.4.\deleted{\unsure{a}} Il est plausible qu'un dévissage plus poussé devrait permettre de prouver des résultats \uline{analogues} à ceux de 6.2, mais sans supposer $U$ \uline{régulier}, et pour les faisceaux $\underline{\Psi}_\nu^q$ eux-mêmes, au lieu des sous-faisceaux $(\underline{\Psi}_\nu^q)_{\mathrm{tame}}$. Pour ce dernier point, il semble qu'il faille auparavant \deleted{\unsure{prouver}}\inserted{d'abord \deleted{\unsure{étudier}}\inserted{résoudre} par l'affirmative la question suivante :} \deleted{ceci ;} \deleted{\unsure{qui n'est pas}}

% L18-21
\uline{Sous les conditions de 3.2}, \uline{supposant les composantes irréductibles} \deleted{$\unsure{Z}$} réduites de $X_0$ \uline{lisses sur} $\overline{k}$, \inserted{et les multiplicités de $X_0$ premières à $p$} \deleted{\unsure{et}} \uline{a-t-on} $(\underline{\Psi}_\nu^{q\deleted{\unsure{s}}})_{\mathrm{tame}}=\underline{\Psi}_\nu^q$, \uline{i.e.} $P$ \uline{opère}\inserted{-t-il} \uline{trivialement sur} \deleted{$\unsure{I}$} $\underline{\Psi}_\nu^q$ ? \deleted{\illegible{1}}
% margin insertion written slanted: et les multiplicités de X_0 premières à p

% L22-25
Cela impliquerait bien sûr, \inserted{si $X$ est propre sur $S$,} que $P$ opère trivialement sur les $H^i(U_{\overline\eta},\Lambda_\nu)$, chose qui n'est pas non plus prouvée à l'heure actuelle. \deleted{\unsure{Et}} Pour le cas $i=1$, c'est vrai, \deleted{\unsure{cf.}} comme on verra dans un exposé ultérieur.

% L26-27 at bottom edge
\deleted{Mais même \unsure{pour|si} $\dim X_0=\dim X_\eta=1$, $U=X_\eta$, $q=1$, \unsure{la réponse à la question d'un point de vue local n'est pas connue}.}
